% generated by roam.py handbook export -- не редактировать, править заметки
% 20261004180905-точная_одномерная_оптимизация.org (body 296c8df052df)

\begin{concept}[{Точная одномерная оптимизация}]\label[concept]{note:8EDDABAC-D2BE-4EC9-8297-DBAC52F70D28}
  Точная одномерная оптимизация~--- \hyperref[note:BEEE81AD-13FC-4846-A160-C72E31760BEE]{выбор шага} $\alpha_k = \arg\min_{\alpha > 0} \phi(\alpha)$, где $\phi(\alpha) = f(x_k + \alpha d_k)$: задача сводится к минимизации функции одной переменной на отрезке.

  Идея этих методов заключается в сужении отрезка, где лежит минимум функции.
  Поэтому они подразумевают \hyperref[note:CD52D818-AA25-471C-97C2-5DBE9787CE53]{унимодальность} функции.

  Примеры (по \hyperref[note:950F2C19-C2C7-42CD-9CDF-0B9AC03A3DC8]{порядку метода}: золотое сечение и параболы~--- нулевого, дихотомия по производной~--- первого):

  \begin{itemize}
    \item \hyperref[note:96DED7D9-3918-45C5-9850-E05030941645]{Метод золотого сечения}
    \item \hyperref[note:124BD9BC-084E-43FB-B617-AF8690B33EF8]{Метод дихотомии по производной}
    \item \hyperref[note:4E950058-62DB-4EDA-A71D-98CACBBC41AA]{Метод парабол}
    \item \hyperref[note:1BCFB956-2053-49A8-ADDF-0D85D53F0610]{Метод Брента}
  \end{itemize}
\end{concept}
